\documentclass[11pt]{article}
\usepackage[margin=1in]{geometry}
\usepackage{amsmath,amssymb,amsthm}
\usepackage[hidelinks]{hyperref}
\setlength{\emergencystretch}{2em}
\newtheorem{theorem}{Theorem}
\newtheorem{lemma}[theorem]{Lemma}
\newtheorem{corollary}[theorem]{Corollary}
\newcommand{\F}{\mathbb F}
\newcommand{\mr}{\operatorname{mr}}
\title{Lift dimension bounds and affine actions with a radical}
\author{Internal research working proof}
\date{16 November 2026}
\begin{document}
\maketitle

The full dimension of a signed-lift space gives a stronger meridional
lower bound than the first successful lift alone. For an $S_4$ quotient,
all independent lift classes can also be realized simultaneously as a
transitive binary affine action preserving a rank-two alternating form.
The resulting targets go beyond the sharp finite Coxeter targets in the
existing small-crossing search. A frozen 16-case test below yields bounds
four, below its certified upper bounds five, and resolves no new knot.
The full Meridional Rank Conjecture remains open. These are written
proofs with independent finite controls, without formal verification,
external review, or a novelty claim.

\section{Every lift contributes to a meridional lower bound}

Let $G_K$ be a knot group and fix a meridian-preserving quotient
$\phi:G_K\twoheadrightarrow S_r$, $r\ge3$, where meridians map to
transpositions. Put $E_r=\{x\in\F_2^r:\sum x_i=0\}$ and
$D_r=E_r\rtimes S_r$. Over a transposition $(pq)$ there are two signed
transpositions, $s_{pq}^a=(a(e_p+e_q),(pq))$, $a\in\F_2$.
The set of all signed transpositions is invariant under conjugation in
$D_r$.

For a presentation on meridians $m_1,\ldots,m_s$, write
$\phi(m_i)=(p_iq_i)$. Seeking lifts $m_i\mapsto s_{p_iq_i}^{a_i}$ makes
every relator a homogeneous linear equation over $\F_2$.
Let $Z_\phi\subset\F_2^s$ be its full solution space, and set
$\ell=\dim Z_\phi$. The canonical split lift exists, so the space is
nonempty. Its definition counts actual homomorphisms with specified
projection and signed meridian image, not homomorphisms modulo conjugacy.

\begin{theorem}\label{thm:count}
The dimension $\ell$ is independent of the chosen meridional presentation,
and $\mr(K)\ge\ell$. Thus $\ell$ as large as a certified bridge upper
bound proves MRC for $K$.
\end{theorem}
\begin{proof}
A relator solution is exactly a homomorphism
$\widetilde\phi:G_K\to D_r$ lifting $\phi$ and taking the specified
meridians to signed transpositions. Since all meridians of a knot are
conjugate, conjugation invariance makes this condition intrinsic; imposing
it on one meridian suffices. Hence there are intrinsically $2^\ell$ such
homomorphisms, independently of the presentation used to compute them.

Let $g_1,\ldots,g_w$ be any meridional generating family. Each
$\widetilde\phi(g_j)$ has just two possible values above the fixed
transposition $\phi(g_j)$. Evaluation on these generators is injective
on the set of homomorphisms. Therefore $2^\ell\le2^w$, so $\ell\le w$.
\end{proof}

This is an elementary counting or restricted-cocycle argument. It is
related to general group-cohomological meridional estimates such as
Gao--Gu--Li--Wang, Lemma 6.2 \cite{gglw}, but does not rely on their
Whitehead-double theorem or its computational inputs. It does not assert
that these homomorphisms are all surjective to $D_r$.

The ambient cut subspace
\[
 C_\phi=\{(b_{p_i}+b_{q_i})_i:b\in\F_2^r\}\subset Z_\phi
\]
has dimension $r-1$ because the projected transposition graph is
connected. Thus $\ell\ge r-1$. The earlier test $\ell>r-1$ obtains a
single transitive $D_r$ quotient and the bound $r$; Theorem~\ref{thm:count}
retains the entire dimension even when $\ell>r$.

For example, the connected sum of $t$ trefoils has meridional
presentation $\langle a,b_1,\ldots,b_t\mid ab_ja=b_jab_j\rangle$.
Map $a$ to $(01)$ and every $b_j$ to $(02)$ in $S_3$. Any choices of
binary signs satisfy these relations in $D_3$: each pair product has
order three. Hence $\ell=t+1$, proving $\mr\ge t+1$. For $t\ge3$ this
exceeds the bound three from a single surjective $D_3\cong S_4$ quotient.
Joining the $t$ standard
two-bridge diagrams gives $b\le t+1$, so equality follows. This recovers
known cases rather than a new connected-sum theorem.

\section{All independent S4 lift classes form an affine action}

Identify the four symbols of $S_4$ with $W=\F_2^2$.
Then $S_4=\operatorname{AGL}(W)$. Let
$B_W(x,y)=x_0y_1+x_1y_0$. Every transposition $s_e$ of the four points
has a unique expression
\[
 s_e(x)=x+(B_W(u_e,x)+\epsilon_e)u_e,\qquad u_e\ne0.
\]
The function $\delta_e(x)=B_W(u_e,x)+\epsilon_e$ is the indicator of its
two moving points. The space $M$ of affine functions $W\to\F_2$ is
three-dimensional. As functions on four points it is exactly the
even-sum permutation module $E_4$, generated by these six indicators.

Let $R=\F_2^d$. Consider
\[
 H_d=\operatorname{Aff}(W,R)\rtimes S_4,
 \qquad |H_d|=24\,2^{3d},
\]
acting on $W\oplus R$ by base affine permutations and fibre translations
given by affine functions of the base point. A vector label $a_e\in R$
on an edge $e=\{p,q\}$ defines
\[
 T_{e,a_e}(x,z)=(s_ex,z+\delta_e(x)a_e).
\]
Define $B((x,z),(y,w))=B_W(x,y)$. It has rank two and radical $R$.
Each generator above is precisely
\[
 T_{e,a_e}(X)=X+(B((u_e,a_e),X)+\epsilon_e)(u_e,a_e).    \tag{1}
\]
Consequently these are affine transvections with the common nonzero
alternating form required by the degenerate extension of the quadratic
bound in \href{PROOF.tex}{the preceding proof}.

\begin{lemma}\label{lem:vectorgraph}
Suppose the underlying edge multigraph on the four base points is
connected. The group generated by the $T_{e,a_e}$ is transitive on
$W\oplus R$ if and only if its cycle-label sums span $R$.
In that case it is all of $H_d$.
\end{lemma}
\begin{proof}
For any function $b:W\to R$, the permutation
$(x,z)\mapsto(x,z+b(x))$ changes edge labels to
$a_e+b(p)+b(q)$. Choose $b$ to make all spanning-tree labels zero.
The corresponding tree transpositions then generate the entire base
$S_4$. Each remaining edge, multiplied by the matching base
transposition, supplies the function $\delta_e a_e$ as a fibre
translation. Conjugating by $S_4$ and adding gives all of $M\otimes U$,
where $U$ is the span of the remaining edge labels, equivalently of all
cycle sums. Thus the generated group is
$(M\otimes U)\rtimes S_4$ after this change of fibre coordinates.

Constant $U$-valued functions belong to $M\otimes U$, so every orbit has
size $4\,2^{\dim U}$. It is transitive precisely when $U=R$.
Even a nonaffine $b$ normalizes $H_d$: the difference
$b(gx)+b(x)$ has even sum on four points, hence is an affine function
in each coordinate. Therefore the conclusion that the original group
is all of $H_d$ also follows when $U=R$.
\end{proof}

\begin{theorem}\label{thm:tower}
For a specified $S_4$ quotient with signed-lift dimension $\ell$, there
is an explicit transitive representation on $2^{\ell-1}$ points,
preserving a rank-two alternating form on a space of dimension $\ell-1$.
Its image is $H_{\ell-3}$ and its meridians have the form (1).
The quadratic bound therefore gives $\mr(K)\ge\ell$ again.
\end{theorem}
\begin{proof}
Choose $d=\ell-3$ vectors $a^{(1)},\ldots,a^{(d)}\in Z_\phi$ whose
classes form a basis of $Z_\phi/C_\phi$. Label edge $i$ by the vector
$(a_i^{(1)},\ldots,a_i^{(d)})$. Each coordinate satisfies the full
relator system, so these labels define a homomorphism into $H_d$.

Over $\F_2$, the cut space is the orthogonal complement of the graph
cycle space, also for parallel edges. Independence of the $a^{(j)}$
modulo cuts means that the cycle-sum map to $R$ has rank $d$.
Lemma~\ref{lem:vectorgraph} now gives the full image and transitivity.
The dimension of $W\oplus R$ is $2+d=\ell-1$; apply the quadratic bound.
\end{proof}

The first images have orders $24,192,1536$ for $\ell=3,4,5$.
The first two are $S_4$ and $W(D_4)$. At $\ell=5$, the image has no
factor five, so cannot surject onto either rank-five target
$W(A_5)=S_6$ or $W(D_5)$ of orders $720$ and $1920$.
The rank-two form is essential to this enlargement in dimension four.

For $P(2,3,\ldots,3)$ on $r$ pretzel bands, use the published cycle
quotient from Baader--Blair--Kjuchukova, Section 2 \cite{bbk}, as before.
Map its first three generators to $(23),(01),(13)$ in $S_4$ and
alternate $(03),(13)$ for the remaining generators. Consecutive pairs
have product order three except the first pair, which is disjoint and
has product order two. The same is true for all scalar signed lifts.
Hence every sign is free and $\ell=r$. The tower gives a transitive
$2^{r-1}$-point action of order $24\,2^{3(r-3)}$, proving the known
equality $\mr=b=r$ in a different family of finite images.

\section{Which small-dimensional affine searches are already covered}

\begin{theorem}\label{thm:small}
For knot groups, transitive common-form affine representations in the
following dimensions force the listed reflection quotients:
\begin{enumerate}
\item dimension two: $W(A_3)=S_4$;
\item dimension three with nonzero alternating form: $W(D_4)$;
\item dimension four with nondegenerate alternating form:
$W(A_5)=S_6$ or $W(D_5)$.
\end{enumerate}
\end{theorem}
\begin{proof}
In dimension two, the affine transvections are the six transpositions
of four points. A transitive group generated by transpositions is $S_4$.

In dimension three, identify the form with the dot product on $E_4$.
Its radical is the all-one vector. The linear form-preserving group
has order $6\cdot4=24$ and is the faithful coordinate $S_4$.
Its six nonradical direction vectors are $e_p+e_q$, so its twelve
affine transvections are precisely the signed transpositions of $D_4$.
A radical direction would give a translation or identity linear part.
Since knot meridians are all conjugate, this would make every meridian
image a translation, hence the image abelian of order at most two,
incompatible with transitivity on eight points.
Disconnected coordinate graphs preserve their component parities on
$E_4$, so are not transitive. The connected signed-graph criterion
from the preceding proof therefore forces the full $D_4$.

For dimension four use the nondegenerate module
$E_6/\langle(1,\ldots,1)\rangle$. Coordinate permutations give a faithful
$S_6$ inside $\operatorname{Sp}(4,2)$; both have order $720$, obtained
for the latter by counting symplectic bases. The fifteen transpositions
are exactly its fifteen linear transvections, with directions
$[e_p+e_q]$. The projected group is generated by transpositions.
All are conjugate within it, so they lie in one connected component on
$m$ symbols and generate $S_m$, with all remaining symbols isolated.
Their direction span has dimension $m-1$ for $m\le5$, and four for
$m=6$. Transitivity of the affine group forces that span to be the
whole four-space, hence $m=5$ or $6$.

If $m=6$, the linear projection is the required $S_6$ quotient. If
$m=5$, realize the module as the even-sum space on these five symbols.
The affine image then lies in $E_5\rtimes S_5$, and its meridians are
signed transpositions. Transitivity forces the full $D_5$ by the same
signed-graph criterion.
\end{proof}

The single meridian conjugacy class is used in the last argument;
it must not be omitted for general links. A rank-two form in dimension
four is also outside that nondegenerate conclusion, as $H_2$ demonstrates.

Section 3.1.1 of the paper by Blair, Kjuchukova and Morrison searches
sharp finite targets
$A_3,A_4,A_5,D_4,D_5,H_3,H_4$ on diagrams through 16 crossings
\cite{bkm}. Thus merely replacing their $D_4,D_5$ search by the earlier
single-lift test cannot add sharp cases on the same diagrams.
The full-dimension test and the towers above are a distinct extension.
This uses their stated search coverage, not a new replay of millions
of diagrams or of all their generating-set lists.

\section{A frozen test and its independent replay}

The author repository was pinned at commit
\texttt{74fe52e57de6f91988f63ac105cbc157534ad966} \cite{code}.
The six downloaded files have checked Git blob identities and SHA256
seals. Its original spreadsheet is read-only; no workbook values or
formula caches were modified. The first sixteen source-order entries
with recorded Wirtinger number five and no sharp finite Coxeter hit
were frozen before testing. This selection does not establish that their
MRC status is otherwise unresolved.

For each diagram a five-seed propagation certificate is produced.
All assignments into transpositions of $S_3$ and $S_4$ are tested,
normalizing the first seed to $(01)$: respectively $81$ and $1296$
assignments per case. Repeated seed labels are allowed, which is
essential for quotients smaller than the upper-bound rank. Every full
colouring and every relator is checked before computing a lift dimension.
The caps are 5000 seed trials and 15 seconds per case, with no replacements.
Every case completed within those caps.

All sixteen have the same normalized distribution: for $S_3$, six
surjective colourings have $\ell=3$ and two have $\ell=4$; for $S_4$,
twelve have $\ell=3$ and twelve have $\ell=4$. Thus the test proves a
lower bound four and retains the certified upper bound five throughout.
It gives no new sharp case. Names, exact rows and all best certificates
are retained in \texttt{lift\_panel.json} and its result file.

An independent verifier reconstructs crossing incidence from cyclic
strand segments, replays seed propagation, enumerates base colourings
with full permutations, and computes each lift dimension by dense
elimination after evaluating one sign variable at a time. It matches
all distributions. Separate controls check the tower relations on every
sheet, group orders $24,192,1536$, connected sums, relabellings and a
redundant conjugate meridian. These are finite checks of the
implementation; Theorems~\ref{thm:count}--\ref{thm:small} are the uniform
mathematical arguments.

The next gate is either an $S_4$ quotient with $\ell\ge5$ on a genuinely
unsettled five-bridge candidate, or a different coefficient quotient.
Repeating small nondegenerate affine searches, or treating failure on
this panel as failure for other targets, would not address the remaining
goal. Original-source audits and non-affine approaches remain part of
the full programme.

\begin{thebibliography}{9}
\bibitem{bbk} S. Baader, R. Blair, A. Kjuchukova,
Coxeter groups and meridional rank of links (2021), Section 2.
\url{https://doi.org/10.1007/s00208-020-02124-z}.
\bibitem{gglw} H. Gao, S. Gu, Z. Li, J. Wang,
Meridional ranks of iterated Whitehead doubles (2026), Section 6.2,
Lemma 6.2, preprint. \url{https://arxiv.org/abs/2611.15673v1}.
\bibitem{bkm} R. Blair, A. Kjuchukova, N. Morrison,
Coxeter quotients of knot groups through 16 crossings, Section 3.1.1.
\url{https://arxiv.org/abs/2208.09032}.
\bibitem{code} N. Morrison, Wirt\_Hm author repository, pinned commit.
\url{https://github.com/ThisSentenceIsALie/Wirt_Hm/tree/74fe52e57de6f91988f63ac105cbc157534ad966}.
\end{thebibliography}
\end{document}
